% TODO make hw stylesheet
\documentclass[11pt]{scrartcl}
\usepackage[a4paper, total={6in, 8in}]{geometry}
\usepackage[usenames,dvipsnames,svgnames]{xcolor}
\usepackage[shortlabels]{enumitem}
\usepackage[framemethod=TikZ]{mdframed}
\usepackage{amsmath,amssymb,amsthm}
\usepackage[colorlinks]{hyperref}

\hypersetup{
	colorlinks=true,
	linkcolor=blue,
	filecolor=magenta,      
	urlcolor=magenta,
	pdftitle={Hough orz}
}

\usepackage{microtype}
\usepackage{mathtools}
\usepackage[headsepline]{scrlayer-scrpage}
\usepackage{thmtools}
\usepackage{todonotes}

\addtolength{\textheight}{3.14cm}
\providecommand{\ol}{\overline}
\providecommand{\eps}{\varepsilon}
\providecommand{\half}{\frac{1}{2}}
\providecommand{\dang}{\measuredangle} %% Directed angle
\providecommand{\CC}{\mathbb C}
\providecommand{\FF}{\mathbb F}
\providecommand{\NN}{\mathbb N}
\providecommand{\QQ}{\mathbb Q}
\providecommand{\RR}{\mathbb R}
\providecommand{\ZZ}{\mathbb Z}
\providecommand{\dg}{^\circ}
\providecommand{\ii}{\item}
\providecommand{\alert}[1]{{\sffamily\textbf{\textcolor{blue}{#1}}}}
\providecommand{\vect}[1]{\overrightarrow{#1}}
\providecommand{\norm}[1]{\left\|#1\right\|}

\reversemarginpar

\mdfdefinestyle{mdbluebox}{roundcorner=10pt,innerbottommargin=9pt,
	linecolor=blue,backgroundcolor=TealBlue!5,}
\declaretheoremstyle[headfont=\sffamily\bfseries\color{MidnightBlue},
mdframed={style=mdbluebox},]{thmbluebox}

\mdfdefinestyle{mdbloobox}{frametitlefont=\bfseries,innerbottommargin=8pt,
	nobreak=true,backgroundcolor=TealBlue!5,linecolor=blue,}
\declaretheoremstyle[headfont=\bfseries\color{MidnightBlue},
mdframed={style=mdbloobox},headpunct={\\[3pt]},postheadspace=0pt,]{thmbloobox}

\mdfdefinestyle{mdredbox}{frametitlefont=\bfseries,innerbottommargin=8pt,
	nobreak=true,backgroundcolor=Salmon!5,linecolor=RawSienna,}
\declaretheoremstyle[headfont=\bfseries\color{RawSienna},
mdframed={style=mdredbox},headpunct={\\[3pt]},postheadspace=0pt,]{thmredbox}

\mdfdefinestyle{mdgreenbox}{linecolor=ForestGreen,backgroundcolor=ForestGreen!5,
	linewidth=2pt,rightline=false,leftline=true,topline=false,bottomline=false,}
\declaretheoremstyle[headfont=\bfseries\sffamily\color{ForestGreen!70!black},
mdframed={style=mdgreenbox},headpunct={ --- },]{thmgreenbox}

\mdfdefinestyle{mdorangebox}{linecolor=RedOrange,backgroundcolor=RedOrange!5,
	linewidth=2pt,rightline=false,leftline=true,topline=false,bottomline=false,}
\declaretheoremstyle[headfont=\bfseries\sffamily\color{RedOrange!70!black},
mdframed={style=mdorangebox},headpunct={ --- },]{thmorangebox}

\mdfdefinestyle{mdblackbox}{linecolor=black,backgroundcolor=RedViolet!5!gray!5,
	linewidth=3pt,rightline=false,leftline=true,topline=false,bottomline=false,}
\declaretheoremstyle[mdframed={style=mdblackbox}]{thmblackbox}

\declaretheoremstyle[spaceabove=6pt,spacebelow=6pt]{basehead}

\declaretheorem[style=thmredbox,name=Problem,numberwithin=section]{problem}
\declaretheorem[style=thmredbox,name=Problem,numbered=no]{problem*}
\declaretheorem[style=thmbloobox,name=Setup]{setup} % for config geo
\declaretheorem[style=thmredbox,name=Example,sibling=problem]{example}
\declaretheorem[style=thmredbox,name=$\bigstar$ Problem,sibling=problem]{niceprob}
\declaretheorem[style=thmbluebox,name=Theorem,sibling=problem]{theorem}
\declaretheorem[style=thmbluebox,name=Corollary,sibling=theorem]{corollary}
\declaretheorem[style=thmbluebox,name=Proposition,sibling=theorem]{prop}
\declaretheorem[style=thmbluebox,name=Lemma,numberwithin=problem]{lemma}
\declaretheorem[style=thmbluebox,name=Theorem,numbered=no]{theorem*}
\declaretheorem[style=thmbluebox,name=Lemma,numbered=no]{lemma*}
\declaretheorem[style=thmgreenbox,name=Claim,numberwithin=problem]{claim}
\declaretheorem[style=thmgreenbox,name=Claim,numbered=no]{claim*}
\declaretheorem[style=thmblackbox,name=Remark,numberwithin=problem]{remark}
\declaretheorem[style=thmblackbox,name=Remark,numbered=no]{remark*}
\declaretheorem[style=thmgreenbox,name=Definition,sibling=theorem]{definition}
\declaretheorem[style=thmgreenbox,name=Definition,numbered=no]{definition*}
\declaretheorem[style=thmorangebox,name=Warning,sibling=theorem]{warning}
\declaretheorem[style=thmorangebox,name=Warning,numbered=no]{warning*}
\declaretheorem[style=thmorangebox,name=Note,numbered=no]{note}
\declaretheorem[style=thmblackbox,name=Exercise,sibling=problem]{exercise}
\declaretheorem[style=thmblackbox,name=Exercise,numbered=no]{exercise*}
\declaretheorem[style=thmblackbox,name=Question,sibling=problem]{question}
\declaretheorem[style=thmblackbox,name=Question,numbered=no]{question*}

\newenvironment{soln}{\begin{proof}[Solution]}{\end{proof}}
\newenvironment{walkthrough}{\noindent \textbf{Walkthrough.}}{\medskip}
\newlist{walk}{enumerate}{3}
\setlist[walk]{label=\bfseries (\alph*)}

\providecommand{\pow}{\operatorname{Pow}}
\providecommand{\vocab}[1]{{\textbf{\textcolor{ForestGreen}{#1}}}}
\providecommand{\lcm}{\operatorname{lcm}}

\usepackage{asymptote}
\begin{asydef}
	defaultpen(fontsize(10pt));
	unitsize(1cm);
	size(8cm); // set a reasonable default
	usepackage("amsmath");
	usepackage("amssymb");
	settings.tex="pdflatex";
	settings.outformat="pdf";
	// Replacement for olympiad+cse5 which is not standard
	import geometry;
	// recalibrate fill and filldraw for conics
	void filldraw(picture pic = currentpicture, conic g, pen fillpen=defaultpen, pen drawpen=defaultpen)
	{ filldraw(pic, (path) g, fillpen, drawpen); }
	void fill(picture pic = currentpicture, conic g, pen p=defaultpen)
	{ filldraw(pic, (path) g, p); }
	// some geometry
	pair foot(pair P, pair A, pair B) { return foot(triangle(A,B,P).VC); }
	pair orthocenter(pair A, pair B, pair C) { return orthocentercenter(A,B,C); }
	pair centroid(pair A, pair B, pair C) { return (A+B+C)/3; }
	// cse5 abbreviations
	path CP(pair P, pair A) { return circle(P, abs(A-P)); }
	path CR(pair P, real r) { return circle(P, r); }
	pair IP(path p, path q) { return intersectionpoints(p,q)[0]; }
	pair OP(path p, path q) { return intersectionpoints(p,q)[1]; }
	path Line(pair A, pair B, real a=0.6, real b=a) { return (a*(A-B)+A)--(b*(B-A)+B); }
	// cse5 more useful functions
	picture CC() {
		picture p=rotate(0)*currentpicture;
		currentpicture.erase();
		return p;
	}
	pair MP(Label s, pair A, pair B = plain.S, pen p = defaultpen) {
		Label L = s;
		L.s = "$"+s.s+"$";
		label(L, A, B, p);
		return A;
	}
	pair Drawing(Label s = "", pair A, pair B = plain.S, pen p = defaultpen) {
		dot(MP(s, A, B, p), p);
		return A;
	}
	path Drawing(path g, pen p = defaultpen, arrowbar ar = None) {
		draw(g, p, ar);
		return g;
	}
\end{asydef}

\usepackage{mathrsfs}

\ihead{\footnotesize Putnam Seminar Solutions}
\ohead{\footnotesize \today}

\title{\vspace{-2em}the grind never stops}
\author{\large Aditya Pahuja}
\date{\today}
\begin{document}
\maketitle
\tableofcontents
\pagebreak

\section{Induction and Pigeonhole}

\begin{problem}
	Prove for all positive integers $n$ the identity
	\[\frac 1{n+1} + \frac 1{n+2} + \cdots + \frac 1{2n} =
	1 - \half + \frac13 - \frac14 - \cdots + \frac1{2n-1} - \frac 1{2n}.\]
\end{problem}
\begin{soln}
	We induct on $n$.
	The base case $n = 1$ is obvious, as
	\[\frac 1{1+1} = \frac 11 - \frac 12 = \half.\]
	
	Suppose now that
	\[\frac 1{n+1} + \frac 1{n+2} + \cdots + \frac 1{2n} =
	1 - \half + \frac13 - \frac14 - \cdots + \frac1{2n-1} - \frac 1{2n}.\]
	Then, add $\frac{1}{2n + 1} + \frac{1}{2n + 2} - \frac{1}{n + 1}$
	to both sides, getting the equality
	\begin{align*}
		\frac 1{n+2} + \cdots + \frac{1}{2n + 1} + \frac 1{2n + 2} &=
		1 - \half + \frac13 - \frac14 - \cdots  - \frac 1{2n}
		+ \frac{1}{2n + 1} + \left(\frac{1}{2n + 2} - \frac{1}{n + 1}\right)\\
		&= 1 - \half + \frac13 - \frac14 - \cdots - \frac 1{2n}
		+ \frac{1}{2n + 1} + \frac{1}{2n + 2} - \frac{2}{2n + 2}\\
		&= 1 - \half + \frac13 - \frac14 - \cdots - \frac 1{2n}
		+ \frac{1}{2n + 1} - \frac{1}{2n + 2}
	\end{align*}
	which is exactly the desired equality with $n + 1$ in place of $n$,
	so we are done by the principle of mathematical induction.
\end{soln}

\begin{problem}
	Let $n$ be a positive integer. Show that
	\[1 + \frac1{2^3} + \frac1{3^3} + \cdots + \frac 1{n^3} < \frac32.\]
\end{problem}

All of the following solutions rely on the idea that
the left-hand side can be bounded above by the obviously convergent series
\[1 + \frac{1}{2^3} + \frac{1}{3^3} + \cdots = \sum_{k=1}^\infty \frac1{n^3}.\]

\begin{proof}[Solution 1]
	Since $\frac1{x^3}$ is decreasing, the integral
	\[1 + \int_1^\infty \frac{1}{x^3}\,dx = \frac32\]
	is an upper bound for the series because the series is just
	the right Riemann sum for the integral.
\end{proof}

\begin{proof}[Solution 2]
	Let $S$ be the sum of all the terms with even denominator,
	so that the series we are trying to bound is $8S$.
	Then, by bounding $\frac{1}{(2n)^3} > \frac{1}{(2n + 1)^3}$, we can say
	that the desired series is bounded above by
	\[8S < 1 + 2S.\]
	This tells us $S < \frac 16$, so $8S < \frac43 < \frac32$.
\end{proof}

\begin{proof}[Solution 3]
	Similar to the previous solution, bound the series by
	\[1 + \frac{2}{2^3} + \frac{4}{4^3} + \frac{8}{8^3} + \cdots,\]
	where each $\frac1{k^3}$ is overestimated by $\frac1{(2^m)^3}$
	where $2^m$ is the largest power of $2$ that is at most $k$.
	This is then a geometric series summing to $\frac43 < \frac32$,
	so we have the necessary upper bound.
\end{proof}

\begin{proof}[Solution 4]
	For $k > 1$, bound
	\[\frac{1}{k^3} < \frac{1}{k^3 - k} =
	\frac{1}{2(k - 1)} + \frac{1}{2(k + 1)} - \frac{1}{k}.\]
	The sum
	\[1 + \sum_{k=2}^\infty \frac{1}{k^3 - k},\]
	which overestimates the desired series, thus telescopes to
	$\frac54 < \frac32$, so we are done.
\end{proof}

\begin{proof}[Solution 5]
	Bound above by
	\[1 + \frac1{2^3} + \frac1{3^3} + \sum_{k=4}^\infty \frac1{k^2}.\]
	It is well known that the last addend is equal to
	$\frac{\pi^2}{6} - 1 - \frac14 - \frac19$, and we can manually check
	(e.g. by bounding $\pi^2 < 10$)
	that this upper bound is less than $\frac32$.
\end{proof}

\begin{problem}
	Let $x_1$, $x_2$, \dots, $x_n$, $y_1$, $y_2$, \dots, $y_m$
	be positive integers, with $m, n>1$. Assume that
	$x_1 + \cdots + x_n = y_1 + \cdots + y_m < mn$.
	Prove that in the expression
	\[x_1 + \cdots + x_n = y_1 + \cdots + y_m\]
	you can suppress some but not all of the terms and still have equality.
\end{problem}

\begin{soln}
	Suppose that $x_1 \le x_2 \le \cdots \le x_n$ and
	$y_1 \le y_2 \le \cdots \le y_m$, and that $x_i \ne y_j$
	for all $i$, $j$ (otherwise we are done by deleting an equal pair of terms),
	and let $S = \sum x_i = \sum y_i$.
	
	We induct on $m$ and $n$ simulatneously,
	showing that the if the statement is true when $(m, n) = (a-1, b)$
	and $(a, b-1)$ then it is also true when $(m, n) = (a, b)$.
	
	First, as the base case, we show that the statement holds
	whenever $m = 2$ (symmetry says that the same logic works for $n=2$).
	To do so, we induct on $n$, with the base case $n = 2$ being trivial.
	We consider two cases:
	\begin{itemize}
		\ii Suppose that some $x_a$ satisfies $2\le x_a \le x_2$.
		Then, delete $x_a$ and replace $y_2$ with $y_2 - x_a$.
		Clearly the new set of $x_i$s has the same sum as the new $y_i$s,
		which is strictly bounded above by $2n - 2 = 2(n - 1)$,
		and we have reduced the number of $x_i$s, so induction finishes.
		\ii Otherwise, every $x_i$ is either larger than $x_2$
		or smaller than $1$. Since $x_2 \ge \frac S2$,
		only $x_n$ can be larger than $y_2$; all the others must be
		equal to $1$, and indeed this forces $x_n > y_2 \ge \frac S2$.
		However, this situation is then impossible
		because $\sum x_i = x_n + n-1$ and
		$y_2 \ge \lceil\frac{x_n + n-1}{2}\rceil$,
		but the right-hand side is always at least as large as $x_n$
		because $x_n\le n$.
	\end{itemize}

	Now, the rest of the problem is simple:
	either $x_n < y_m$, in which case we delete $x_n$
	and replace $y_m$ with $y_m - x_n$ to reduce to $(m, n-1)$,
	or we reduce to $(m, n-1)$ in the same way.
\end{soln}

\begin{problem}
	no.
\end{problem}

\begin{problem}
	The vertices of a convex polygon are colored by at least three colors
	so that no two consecutive vertices have the same color.
	Prove that you can dissect the polygon into triangles by diagonals
	that do not cross and whose endpoints have different colors.
\end{problem}

\begin{soln}
	Very simple induction, actually --- we only need to be able to
	always chop a triangle off. Supposing the polygon has $n$ sides,
	the base case $n=3$ is trivial.
	
	Then, consider an $n$-gon $P_1P_2\cdots P_n$.
	If there exists a color that is only assigned to one vertex $P_a$,
	then just draw the edge $P_aP_{a+2}$ and look at the $n$-gon
	obtained by deleting $P_{a+1}$, which is still convex
	and still has at least three color; namely, $P_{a-1}$
	and $P_{a-2}$ both differ in color from each other and from $P_a$,
	since $P_a$ was assumed to be the only vertex of its color.
	The inductive hypothesis then applies, and we are done.
	
	Otherwise, consider the pairs
	$(P_i, P_{i+2})$ with indices mod $n$. Suppose FTSOC each pair
	is monochromatic; then, if $n$ is odd, all the vertices are the same color,
	or, if $n$ is even, then there are at most two colors in use,
	both of which violate the assumptions about the polygon.
	Thus, there exists a pair $(P_a, P_{a+2})$ with differing colors;
	draw this edge and delete $P_{a+1}$. By assumption,
	$P_{a+1}$ is not the only vertex of its color, so we have not
	decreased the number of available colors,
	and thus the inductive hypothesis applies here, too.
\end{soln}

\begin{problem}
	A sequence of $m$ positive integers contains exactly $n$ distinct terms.
	Prove that if $2^n\le m$ then there exists a block of consecutive terms
	whose product is a perfect square.
\end{problem}

\begin{soln}
	Let $x_1$, \dots, $x_n$ be the possible $n$ distinct values in the sequence.
	Then, replace the $k$th term in the sequence with the tuple
	$(0, 0, \dots, 0, 1, 0, \dots, 0)$ whose $i$th entry is $1$
	if and only if the the term in the sequence is $x_i$.
	Treat these tuples as vectors in $\FF_2^n$
	and let $v_1$, $v_2$, \dots, $v_m$ be the sequence of vectors.
	
	Now, consider the sequence of partial sums
	$v_1$, $v_1 + v_2$, $v_1 + v_2 + v_3$, \dots, $v_1 + \cdots + v_m$.
	If this sequence has two equal values, then we are done:
	specifically, if $i < j$,
	\[v_1 + \cdots + v_i = v_1 + \cdots + v_j \implies
	v_{i+1} + \cdots + v_j = (0, 0, \dots, 0),\]
	meaning that the product of the terms in the sequence corresponding to the
	$v_k$ between $v_{i+1}$ and $v_j$ is a square.
	More explicitly, the product is the product of $x_i^{e_i}$
	over all $i$, where $e_i$ is the number of times
	that $1$ appears in the $i$th entry within our range of terms ---
	since our vectors are in $\FF_2^n$, the sum being the zero vector
	means that the $e_i$ are all even, hence square.
	Similarly, if it contains the zero vector, we are also done.
	There are only $2^n - 1$ nonzero vectors,
	so pigeonhole tells us that two of the $m\ge 2^n$ vectors are
	equal and nonzero if no zero vectors are in the partial sums,
	which is what we wanted.
\end{soln}

\begin{problem}
	Let $p$ be a prime number and $a$, $b$, $c$ be in $\FF_p$ such that
	$a$ and $b$ are nonzero. Prove that there exist $x,y\in\FF_p$ such that
	$ax^2 + by^2 = c$.
\end{problem}

\begin{soln}
	If $c = 0$ then $x = y = 0$ works. Suppose $c\ne 0$.
	
	If $c$ isn't a QR mod $p$, multiply the equation by some (nonzero) non-QR
	to get the equivalent equation
	\[a'x^2 + b'y^2 = c'.\]
	Now, if one of $a'$ or $b'$ is a QR (WLOG $a'$) then
	$a'x^2$ is free to be any QR, which means we can set
	$a'x^2 = d^2$ and $y = 0$.
	
	The other case is where $a'$ and $b'$ are not QRs.
	In this case we look back at the original equation,
	where $a$ and $b$ are QRs and $c$ isn't.
	Then, $ax^2$ and $by^2$ are arbitrary QRs.
	By Fermat's two squares theorem, $u^2 + v^2$ can be any prime number
	that is $1\pmod 4$, and by Dirichlet's theorem,
	there are infinitely many such primes in every nonzero residue class
	modulo $p$. Pick a prime $q$ that is congruent
	to both $c\pmod p$ and $1\pmod 4$
	(which exists either by CRT or, if $p = 2$, trivially, since we
	already dealt with $c = 0$),
	apply the theorem for $ax^2 = u^2$ and $by^2 = v^2$
	(equality in $\FF_p$) and win.
\end{soln}

\begin{problem}
	Show that there is a positive term of the Fibonacci sequence
	that is divisible by $1000$.
\end{problem}

\begin{soln}
	We prove the well-known fact that the Fibonacci sequence is periodic.
	First, by Pigeonhole on the pairs $(F_i, F_{i+1})$ modulo $1000$,
	some pair has to appear twice, i.e. $(F_n, F_{n+1})\equiv(F_{m},F_{m+1})$
	modulo $1000$ for some $m > n$.
	Then, we can backfill the numbers before $F_n$ to find that
	the sequence is purely periodic (rather than just eventually),
	which means some nonzero $F_n$ that is congruent to $F_0 = 0$ mod $1000$.
\end{soln}

\begin{problem}
	Let $m$ be a positive integer.
	Prove that among any $2m + 1$ distinct integers
	of absolute value less than $2m - 1$,
	there exist three whose sum is zero.
\end{problem}

\begin{soln}
	If one of the integers is zero, we're done
	because one of the pairs $(1, -1)$, $(2, -2)$, \dots, $(2m-2, -2m+2)$
	must also be selected.
	
	Now, suppose none of them are zero. Then, say there are $k$
	positive integers in the set, $a_1 < a_2 < \cdots < a_k$,
	and clearly $k\ge 3$.
	Consider the sums
	\[a_1 + a_2 < a_1 + a_3 < a_2 + a_3 < a_2 + a_4 < a_3 + a_4 < \cdots.\]
	There are $2k-3$ of these.
	We want to show that one of these sums is equal in absolute value
	to one of the chosen negative numbers.
	Since there are $2m - 2$ choosable negative numbers
	and $2m+1-k$ chosen negative numbers, $k-3$ negative numbers are not chosen.
	However, $2k - 3 > k - 3$, meaning one of the sums must be equal
	in absolute value to one of the chosen negative numbers, as desired.
\end{soln}

\begin{remark*}
	Absolute value \emph{at most} $2m-1$ also works.
\end{remark*}

\begin{problem}
	Let $P_1$, \dots, $P_{2m}$ be a permutation of the vertices of a
	regular polygon. Prove that the closed polygonal line $P_1P_2\dots P_{2m}$
	contains a pair of parallel segments.
\end{problem}

\begin{soln}
	Let $\omega = \exp(\frac{\pi i}{m})$, so that the vertices of the
	$2m$-gon are powers of $\omega$. Then, if $P_j = \omega^{x_j}$
	for $1\le x_j\le 2m$ and for all $j$,
	consider the permutation $(x_1, \dots, x_{2m})$ of $(1, 2, \dots, 2x)$.
	The main idea is the following fact
	about complex numbers on the unit circle:
	\begin{claim}
		If $a$, $b$, $c$, and $d$ lie on the unit circle centered at $0$
		in the complex plane,
		then the segment joining $a$ and $b$ is parallel to
		the segment joining $c$ and $d$ if and only if $ab = cd$.
	\end{claim}
	\begin{proof}
		The segments being parallel means that the expression
		\[\frac{b - a}{d - c}\]
		is real, since the numerator can be thought of as
		a vector along the line through $a$, $b$ and the denominator as
		a vector along the line through $c$, $d$.
		This means the expression is equal to its own conjugate
		($z = \ol z$ if and only if $z\in\RR$), so
		\[\frac{b-a}{d-c} = \ol{\left(\frac{b-a}{d-c}\right)}
		= \frac{\ol b-\ol a}{\ol d- \ol c}
		= \frac{\frac1b - \frac1a}{\frac 1d - \frac 1c}\]
		because $|z| = 1$ means $z\ol z = 1$, i.e. $\ol z = \frac 1z$.
		This simplifies to $ab = cd$ as desired.
	\end{proof}
	Since $\omega^{x_a}\cdot \omega^{x^b} = \omega^{x_a + x_b}$
	where the exponent is only relevant as far as its residue modulo $2m$,
	we just need to show that there exist $u$ and $v$ such that
	\[x_u + x_{u+1}\equiv x_v + x_{v+1}\pmod{2m},\]
	noting that $x_u$, $x_{u+1}$, $x_{v}$, and $x_{v+1}$ must all be distinct
	because clearly segments sharing a vertex can't be parallel,
	as each of the $P_j$s are distinct.
	Now, if each $x_k + x_{k+1}$ was distinct mod $2m$ so that
	\[\{x_k + x_{k+1} : 1\le k\le 2m\} = \{1, 2, \dots, 2m\},\]
	then the sum of the $x_k + x_{k+1}$ is on one hand
	\[2\cdot(1 + 2 + \cdots + 2m) = 2m(2m + 1)\equiv 0\pmod 2m\]
	because each element of the permutation appears twice,
	but on the other hand is equal to
	\[1 + 2 + \cdots + 2m = m(2m + 1)\equiv m\pmod{2m}\]
	by the distinctness assumption, but this is a contradiction.
	Thus two of the $x_k + x_{k+1}$ are identical,
	which is what we wanted.
\end{soln}

\begin{problem}
	The points of the plane are colored with one of a finite number of colors.
	Prove that some rectangle has all vertices the same color.
\end{problem}

\begin{soln}
	We just focus on $\ZZ^2$. Suppose there are $n$ colors, and
	consider the $n^{n+1} + 1$ sets of points
	\[S_k = \{(i, k) : 1\le i\le n + 1\}\]
	where $1\le k\le n^{n+1}+1$.
	Then, by pigeonhole, some two $S_k$s have identical color patterns;
	that is, there exist $a$, $b$ such that
	$(i, a)$ and $(i, b)$ have the same color for all $1\le i\le n+1$,
	because there are only $n^{n+1}$ possible colorings for any $S_k$.
	To finish, we see that within $S_a$,
	there exist two points of the same color
	because there are $n+1 > n$ points in $S_a$:
	call these  $(c, a)$ and $(d, a)$, and then
	we are done because $(c, a)$, $(d, a)$, $(c, b)$, $(d, b)$
	make a monochromatic rectangle.
\end{soln}

\begin{problem}
	Inside the unit square lie several circles, the sum of whose circumferences
	is at least $10$. Prove that infinitely many lines intersect
	at least four circles.
\end{problem}

\begin{soln}
	Let the unit square have vertices $(0,0)$, $(1,0)$, $(0,1)$, $(1,1)$.
	Then, for each $0 < r < 1$, let $f(r)$ be the number of circles
	which have an interior point with $x$-coordinate $r$.
	The integral of $f$ over $(0,1)$ is then $\frac{10}{\pi} > 3$,
	so there exists a set of $r$-values with nonzero measure
	for which $f(r) = 4$, the end.
\end{soln}

\pagebreak
\section{Invariants}

\begin{problem}
	TODO fill in
\end{problem}

\begin{problem}
	content
\end{problem}



\end{document}